\section{LLM 工作原理速览}

\subsection{基本架构}

LLM（Large Language Model）是一种自回归（autoregressive）的 next-token prediction 模型。其基本工作流程为：

\begin{enumerate}
    \item \textbf{Tokenization}：使用固定词汇表将输入文本分割为 token
    \item \textbf{Embedding}：将 token 转换为固定维度的数值向量（通常 1K--3K 维）
    \item \textbf{Transformer layers}：经过 12--96+ 层 Transformer（使用 self-attention 机制，Vaswani et al. 2017）
    \item \textbf{输出概率分布}：在词汇表上生成下一个 token 的概率分布
\end{enumerate}

\begin{knowledgebox}{Transformer 架构的核心}
Self-attention 机制允许模型在生成每个 token 时``看到''输入序列中的所有其他 token，从而捕获长距离依赖关系。这是 LLM 能够理解代码上下文的关键技术基础。
\end{knowledgebox}

\subsection{训练过程：三阶段范式}

现代 LLM 的训练通常遵循三阶段范式：

\textbf{Stage 1: Self-supervised Pretraining}
\begin{itemize}
    \item 在大量公开数据上进行自监督预训练
    \item 数据源包括 Common Crawl、Wikipedia、StackExchange、公开 GitHub 仓库等
    \item 数据规模：数千亿到万亿+ token（语言和代码）
    \item 目标：学习语言的基本结构和世界知识
\end{itemize}

\textbf{Stage 2: Supervised Fine-tuning (SFT)}
\begin{itemize}
    \item 教模型遵循指令
    \item 使用高质量的 prompt-response 对进行训练
    \item 数据规模：数万到数十万对
    \item 例：``Write a for loop'' $\rightarrow$ ``ok here's a for loop...''
\end{itemize}

\textbf{Stage 3: Preference Tuning (RLHF/DPO)}
\begin{itemize}
    \item 将模型输出与人类偏好对齐（helpfulness、correctness、readability）
    \item 收集同一 prompt 的多个输出，训练 reward model 预测人类偏好
    \item 数据规模：数万到数十万个人类标注的比较对
\end{itemize}

\begin{importantbox}{Reasoning Models 的演进}
在三阶段基础上，reasoning models 进一步扩展了训练：加入 chain-of-thought reasoning traces、工具使用整合、对推理步骤的人类偏好收集，以及通过强化学习（RL）学习如何评估推理轨迹和回溯（backtrack）。
\end{importantbox}

\subsection{模型规模参考}

\begin{itemize}
    \item GPT-3 / Claude 3.5 Sonnet: $\sim$175B parameters
    \item LLaMA 3.1: 405B parameters
    \item GPT-4: $\sim$1.8T parameters（传闻）
\end{itemize}

\subsection{LLM 编程的优势与局限}

\textbf{优势}：
\begin{itemize}
    \item Expert-level 代码补全
    \item 代码理解与解释
    \item 代码修复（bug fixing）
\end{itemize}

\textbf{局限}：
\begin{itemize}
    \item \textbf{Hallucination}：生成不存在或过时的 API（可通过 robust context engineering 缓解）
    \item \textbf{Context window 限制}：$\sim$100--200K tokens，但并非所有 token 都同等有效
    \item \textbf{延迟}：每次请求数秒到数分钟（需要据此规划任务委派）
    \item \textbf{成本}：最佳模型的输入 token \$1--3/M，输出 token \$10+/M
\end{itemize}

\begin{warningbox}{Context Window 的质量问题}
虽然现代 LLM 宣称支持 100K--200K token 的 context window，但研究表明模型对窗口中间部分的信息关注度显著下降（``lost in the middle''现象）。不能简单地把所有内容塞入 context 就期望好结果。
\end{warningbox}

\subsection{本章小结}

LLM 本质上是 next-token prediction 模型，通过三阶段训练（预训练 $\rightarrow$ SFT $\rightarrow$ 偏好对齐）获得强大的代码生成能力。理解其工作原理和局限性是有效使用 AI 编程工具的前提。

